Le nombre $1066$ nous fait immédiatement penser à la bataille d'Hastings et la victoire de Guillaume le Conquérant, duc de Normandie~;
$1776$ à l'indépendance des États-Unis~; $1789$ à la Révolution
française. Quand on lit $3.14159$ on sait qu'il s'agit d'une valeur approchée de $\pi$
(écrit avec la convention anglo-saxonne qui utilise un point à la place de notre virgule).
Par contre quand on lit le chiffre $2$, on ne pense à rien de précis et il est à peu près sûr
qu'on se demande 
\raisebox{0.25ex}{$\scriptscriptstyle \boldmath\ll$}~%
deux quoi$\:$?%
~\raisebox{0.25ex}{$\scriptscriptstyle \boldmath\gg$}.
Tout cela pour dire que lorsqu'on voit un nombre dans un
programme il n'est pas toujours facile de comprendre immédiatement de quoi il s'agit. Pour le savoir,
on peut, en \pascal\ nommer le nombre à l'aide d'une déclaration de constante placée au
début du programme~:

\vspace{\neglskip}\mbox{\blc  \blk\idt  \mcl const\idt }\par 
\vspace{\neglskip}\mbox{\blc  \blk\blk\idt  RevolutionFrancaise\ $=$\ 1789$;$\ \ \cmt\{\ Date de la révolution\ \}}\par 

\noindent puis utiliser le nom {\idt RevolutionFrancaise} à la place du nombre 
{\small\rm 1789} dans la suite du programme.


\begin{exemple}[Croix de Malte]{\rm
Dans le programme de dessin d'une croix de Malte ci-dessous, on
utilise une constante {\idt Longueur} pour définir la longueur d'une
branche (c'est-à-dire la distance de sa base au centre de la croix) et
une autre {\idt Largeur} pour définir la largeur d'une branche
(c'est-à-dire la distance entre les extrémités de sa base)~:\hfill\\[22mm]
\includePostScript{CroixDeMalte_1}{0.95}{0.95}%
\vspace*{-22mm}%
\input CroixDeMalte_1.tex

Pour choisir la meilleure forme possible pour
cette croix, il n'y a qu'un seul endroit dans le programme où il est nécessaire
de modifier ces constantes. Par exemple pour obtenir la plus petite croix représentée à la
figure \ref{fig:crmadiff} ci-dessous, il suffit de prendre~:\par
\vspace{\neglskip}\mbox{\blc  \blk\idt  \mcl const\idt }\par 
\vspace{\neglskip}\mbox{\blc  \blk\blk\idt  Largeur\ $=$\ 10$;$\ \ \cmt  \{\ Largeur\ d{'}une\ branche\ \ \}\idt }\par 
\vspace{\neglskip}\mbox{\blc  \blk\blk\idt  Longueur\ $=$\ 15$;$\ \cmt  \{\ Longueur\ d{'}une\ branche\ \}\idt }\par 
\begin{figure}[htbp]
 \includevspacePostScript[2.5cm]{CroixDeMalteDiff}{0.95}{0.95}
 \caption{Croix de Malte $10\times 15$, $15\times 20$, $20\times 25$, $25\times 30$, 
 $10\times 35$, $35\times 35$}
 \label{fig:crmadiff}
\end{figure}
\noindent $\Box$ }\end{exemple}

On indique que l'on utilise une ou plu\-sieurs {\em constantes nommées}
par des {\em déclarations de constantes}~:\index{Déclaration de constantes}

\begin{center}
\begin{tabular}{@{}|p{12.5cm}|@{}} \hline
{\em Déclaration de constantes~:}\\ \hline \hline
\begin{minipage}[t]{12.5cm}
{\medskip
\mbox{\blc\ \ \ \ {\mcl const}}\par 
\mbox{\blc\ \ \ \ \ \idt {\small\rm Identificateur1} $=$\ {\small\rm Nombre1}$;$}\par 
\mbox{\blc\ \ \ \ \ \idt {\small\rm Identificateur2} $=$\ {\small\rm Nombre2}$;$}\par 
\mbox{\blc\ \ \ \ \ \small\rm  \ldots}\par 
\medskip}
\end{minipage}\\ \hline 
\end{tabular}
\end{center}

\noindent Le mot anglais {\mcl const} (abréviation de `constant') doit 
obligatoirement être sé\-pa\-ré de
l'iden\-ti\-fi\-ca\-teur qui le suit par un blanc \hspace{0.3em}\blanc\hspace{0.3em} ou un retour à la ligne.
Après ces déclarations, on peut utiliser les noms de constantes
`{\small\rm Identificateur1}' à la place de `{\small\rm Nombre1}' et `{\small\rm Identificateur2}'
à la place de `{\small\rm Nombre2}', etc.

Dans un programme \pascal, une déclaration définit des constantes, des variables ou des
procédures~:\index{Déclaration}

\begin{center}
\begin{tabular}{@{}|p{12.5cm}|@{}} \hline
{\em Déclaration~:}\\ \hline \hline
\begin{minipage}[t]{12.5cm}
{\medskip
\mbox{\blc\ \ \ \ \small\rm  Déclaration de constantes}\par 
\mbox{\blc\ \ \tiny\rm  ou}\par 
\mbox{\blc\ \ \ \ \small\rm  Déclaration de variables}\par 
\mbox{\blc\ \ \tiny\rm  ou}\par 
\mbox{\blc\ \ \ \ \small\rm  Déclaration de procédures}\par 
\medskip}
\end{minipage}\\ \hline 
\end{tabular}
\end{center}

Les déclarations de constantes, variables ou procédures se placent, dans n'importe quel
ordre, avant le corps du programme~:\index{Programme}

\begin{center}
\begin{tabular}{@{}|p{12.5cm}|@{}} \hline
{\em Programme \pascal\ :}\\ \hline \hline
\begin{minipage}[t]{12.5cm}
{\medskip
\mbox{\blc\ \ \ {\mcl program}\hspace{0.3em}\blanc\hspace{0.3em}\idt NomDuProgramme$;$}\par 
\mbox{\blc\ \ \ \ \idt {\mcl uses}\ \ldots\ $;$}\par 
\mbox{\blc\ \ \ \ \small\rm  \ldots\ Déclarations du programme \ldots}\par 
\mbox{\blc\ \ \ \mcl begin}\par 
\mbox{\blc\ \ \ \ \small\rm  \ldots\ Instructions du programme\footnote{\ Utilisant les constantes,
variables et procédures déclarées.} (séparées par des `$;$') \ldots}\par 
\mbox{\blc\ \ \ {\mcl end}\idt$.$}\par 
\medskip}
\end{minipage}\\ \hline 
\end{tabular}
\end{center}

La principale règle à respecter est que tout identificateur doit être déclaré avant d'être
utilisé.

\noindent \rule{137mm}{0.1mm} %---------------------------------------------
\begin{exercice}{Écrire des programmes \pascal\ comportant des constantes nom\-mées
pour dessiner~:

\tiret l'hélice singulière de la figure \ref{fig:hesing} ci-dessous$\:$;

\tiret le canardeau de la figure \ref{fig:canar} ci-dessous.

\noindent $\Box$}\end{exercice}
\begin{figure}[htbp]
 \includevspacePostScript[4.6cm]{HeliceSinguliere}{0.95}{0.95}
 \caption{Hélice singulière}
 \label{fig:hesing}
\end{figure}
\begin{figure}[htbp]
 \includevspacePostScript[4.2cm]{Canardeau}{0.95}{0.95}
 \caption{Canardeau sortant de son œuf}
 \label{fig:canar}
\end{figure}
\begin{corrige}{
Nous donnons un programme, l'autre se trouve sur la disquette  d'accompagnement du livre~:
\input HeliceSinguliere.tex
\noindent $\Box$}\end{corrige}
\nopagebreak

%%%%%%%%%%%%%%%%%%%%%%%%%%%%%%%%%%%%%%%%%%%%%%%%%%%%%%%%%%%%%%%%%%%%%%%%%%%
\noindent \framebox[\textwidth]{
\begin{minipage}{\myparboxwidth}
\setlength{\parindent}{2em}
\small \sf
\vspace{\abovedisplayskip}
\section*{Addenda}

Une {\em déclaration de constante} {\mcl const} permet de donner un nom  
à un nombre particulier. Par la suite le nom, plus significatif, peut être 
employé à la place du nombre. L'avantage immédiat est que pour changer de nombre il
suffit de le faire dans la déclaration de constante alors qu'autrement il faut rechercher toutes ses
occurrences  dans le programme.

Ceci peut être comparé à la notion de variable en mathématiques c'est-à-dire à l'emploi 
de lettres $x$, $y$, $z$ \ldots pour désigner des nombres quelconques. La différence est qu'en
\pascal\ le nom de constante désigne une valeur particulière alors qu'en mathématiques il
s'agit d'une valeur quelconque.

\vspace{\abovedisplayskip}
\end{minipage}
}
%%%%%%%%%%%%%%%%%%%%%%%%%%%%%%%%%%%%%%%%%%%%%%%%%%%%%%%%%%%%%%%%%%%%%%%%%%%%
%%%%%%%%%%%%%%%%%%%%%%%%%%%%%%%%%%%%%%%%%%%%%%%%%%%%%%%%%%%%%%%%%%%%%%%%%%%
\noindent \framebox[\textwidth]{
\begin{minipage}{\myparboxwidth}
\setlength{\parindent}{2em}
\small \sf
\vspace{\abovedisplayskip}
Historiquement les premiers problèmes mathématiques traités par les babyloniens (1800 ans av. J.-C.)
comportaient une solution algorithmique donnée pour un exemple et des tables numériques
 séparées pour résoudre d'autres problèmes similaires \cite{ritter1}.
Au Moyen Age on utilisait le langage rhétorique comme `un nombre ajouté à son carré
valent 12' pour `$x + x^{2} = 12$'. Chuquet\footnote{\ \href{https://fr.wikipedia.org/wiki/Nicolas_Chuquet}{Nicolas Chuquet}, mathématicien 
français, v. 1445-1500.}, auteur du plus ancien traité d'algèbre écrit par un 
français (1484), utilisait $1^{1}$ pour désigner l'inconnue, $1^{2}$ pour désigner son
carré et les symboles $\bar{p}$ et $\bar{m}$ pour l'addition et la soustraction. 
Il pouvait écrire `$12\;\bar{p}\;3^{1}$ egaulx a $4^{2}$' pour `$12 + 3x = 4x^{2}$'. 
Viète\footnote{\ \href{https://fr.wikipedia.org/wiki/François_Viète}{François Viète}, mathématicien français, 1530-1603.} fit faire un pas 
important vers la symbolisation en algèbre en introduisant l'usage des lettres pour désigner
les quantités inconnues. Il écrivait vers 1590~: 
`$1\;$ QC $\;-\;15\;$QQ$\;+\;85\;$C$\;+\;225\;$Q$\;+\;274\;$N égalent $120$', soit si
C si\-gni\-fie cube, donc $x^{3}$, Q si\-gni\-fie carré, donc $x^{2}$, et N signifie nombre, donc $x$, l'équation
$x^{5} - 15x^{4} + 85x^{3} + 225x^{2} + 274x = 120$ (\cite{benoit1}). Il fallut at\-ten\-dre qua\-tre mathématiciens 
français
Desargues\footnote{\ \href{https://fr.wikipedia.org/wiki/Girard_Desargues}{Girard Desargues}, mathématicien français, 1591--1661.},
Descartes\footnote{\ \href{https://fr.wikipedia.org/wiki/René_Descartes}{René Descartes}, mathématicien français, 1596--1650.},
Fermat\footnote{\ \href{https://fr.wikipedia.org/wiki/Pierre_de_Fermat}{Pierre de Fermat}, mathématicien français, 1601--1665.} et
Pascal\footnote{\ \href{https://fr.wikipedia.org/wiki/Blaise_Pascal}{Blaise Pascal}, mathématicien français, 1623--1662.} (ce dernier ayant été choisi
par Wirth pour nommer son langage de programmation \pascal) pour entrer dans l'ère
moderne des mathématiques.

\vspace{\abovedisplayskip}
\end{minipage}
}
%%%%%%%%%%%%%%%%%%%%%%%%%%%%%%%%%%%%%%%%%%%%%%%%%%%%%%%%%%%%%%%%%%%%%%%%%%%%
